% ===========================================================================
%  CAPITULO 1 — Como ler este comparativo e premissas inegociaveis
% ===========================================================================
\chapter{Como ler este comparativo e premissas da plataforma}
\label{cap:como-ler}

Este documento analisa, peça por peça, a infraestrutura em nuvem da
\textbf{Plataforma Unooria} e compara a arquitetura atualmente implantada no
\textbf{Google Cloud Platform (GCP)} com as suas traduções arquiteturais diretas
em \textbf{Amazon Web Services (AWS)} e \textbf{Microsoft Azure}.

Diferente de um comparativo genérico de mercado, este estudo parte das
\textbf{invariantes reais do código e do modelo de negócio} da Unooria ---
documentadas nas ADRs do monorepo e nos módulos Terraform do repositório de
fundação. O objetivo é responder a três perguntas práticas:

\begin{enumerate}
  \item \textbf{Equivalência técnica:} Qual serviço de cada nuvem cumpre
    exatamente o papel das peças atuais da Plataforma, das Experiências, da API
    e do Uso de IA sem violar as regras invioláveis de segurança, \emph{tenancy}
    e tempo de resposta?
  \item \textbf{Custo de operação para uma pessoa:} Onde a troca de nuvem reduz
    ou multiplica a carga cognitiva e operacional de um único desenvolvedor
    mantendo múltiplos contratos B2G?
  \item \textbf{Estrutura de custos (TCO) e acoplamento:} Como se comportam o
    piso fixo mensal por contrato e o custo variável de IA nas três nuvens, e
    quais componentes têm portabilidade imediata versus acoplamento profundo?
\end{enumerate}

\section{Estrutura do documento e convenções visuais}

O estudo está organizado nos quatro eixos solicitados, precedidos por uma visão
executiva de equivalência e seguidos pela análise financeira e de migração:

\begin{itemize}
  \item \textbf{Capítulo~\ref{cap:visao-geral} --- Visão geral e matriz de equivalência:}
    diagrama comparativo da topologia completa e tabela sintética de mapeamento
    serviço a serviço (\cd{GCP} $\times$ \cd{AWS} $\times$ \cd{Azure}).
  \item \textbf{Capítulo~\ref{cap:eixo-plataforma} --- Eixo 1: Plataforma:}
    borda única (Load Balancer, TLS e WAF), hospedagem do \arq{apps/hub}
    (Next.js App Router), banco de dados PostgreSQL com \emph{Row-Level Security}
    (RLS), mensageria de telemetria (\emph{push} autenticado), identidade
    federada (\emph{keyless}) e observabilidade (logs, métricas, alertas e OTLP).
  \item \textbf{Capítulo~\ref{cap:eixo-experiencias} --- Eixo 2: Experiências:}
    entrega de \emph{bundles} estáticos imutáveis por versão (\arq{apps/experiencias/*}),
    estratégia de \emph{cache} e \emph{rollback} em segundos via CDN, isolamento
    de origem para \cd{postMessage}, mídia pesada (\cd{-assets}), guarda
    versionada de planos de aula (\cd{-planos}) e \emph{sourcemaps} privados.
  \item \textbf{Capítulo~\ref{cap:eixo-api} --- Eixo 3: API:}
    runtime \emph{serverless} do \arq{apps/api} (NestJS), alocação contínua de
    CPU para \emph{workers} em segundo plano pós-\cd{HTTP 202}, concorrência por
    instância, \emph{timeouts} longos para \emph{streams} SSE e WebSockets de
    voz ao vivo, saída privada de rede (VPC) e execução de \emph{migrations} via
    \emph{job} isolado.
  \item \textbf{Capítulo~\ref{cap:eixo-ia} --- Eixo 4: Uso de Inteligência Artificial:}
    a camada única (\cd{AiProvider} com texto curto, saída longa de 64k
    \emph{tokens}, JSON validado, geração de imagem com proporção nativa e TTS),
    o Agente de IA do Professor (\arq{apps/agente-professor} em Python/ADK com
    sessões gerenciadas e dupla busca híbrida estruturada/layout com
    \cd{pageSpan}) e a Ajuda por voz bidirecional (\cd{Gemini Live}) e vídeo
    (\cd{LiveAvatar}).
  \item \textbf{Capítulo~\ref{cap:custos-migracao} --- Custos, operação solo e matriz de acoplamento:}
    estimativa comparada do piso fixo mensal por contrato e consumo variável de
    IA, impacto na rotina de operação de uma única pessoa e esforço real de
    portabilidade de cada peça.
\end{itemize}

\subsection{Selos de aderência arquitetural}

Cada componente analisado nos Capítulos~\ref{cap:eixo-plataforma}
a~\ref{cap:eixo-ia} traz uma ficha técnica e selos que classificam o grau de
aderência de cada nuvem aos requisitos da Unooria:

\begin{tabularx}{\linewidth}{@{}L{42mm}X@{}}
\toprule
\headrow \thd{Selo} & \thd{Significado técnico neste estudo}\\
\midrule
\nativo &
  O serviço da nuvem atende ao requisito de forma nativa e direta, com
  configuração declarativa em Terraform e sem exigir peças intermediárias nem
  mudanças no código da aplicação.\\
\zebra \equivalente &
  A nuvem oferece capacidade funcional e de segurança equivalente, mas exige a
  combinação de dois recursos (ex.: fila + função consumidora, ou proxy de
  conexões separado) ou pequenos ajustes em variáveis/adaptadores.\\
\adaptacao &
  A nuvem atende ao caso de uso, porém impõe mudança relevante de arquitetura,
  reescrita de cliente/SDK (ex.: refatorar o agente ADK ou o pipeline de
  \emph{chunking} de PDFs com páginas) ou salto expressivo no piso fixo de custo.\\
\zebra \incompativel &
  O serviço candidato viola uma premissa técnica da plataforma (ex.: limite
  rígido de 15 minutos ou congelamento de CPU no AWS Lambda para conexões
  WebSocket de 60 minutos e \emph{jobs} \emph{in-process} pós-202).\\
\bottomrule
\end{tabularx}

\section{Premissas de negócio e arquitetura que a nuvem precisa respeitar}

Qualquer comparação honesta entre GCP, AWS e Azure precisa julgar os serviços
contra as restrições reais do produto. Uma solução que exige um time de
SRE dedicado ou que encarece o piso ocioso de cada contrato em milhares de
dólares é inviável para o modelo da Unooria, ainda que seja tecnicamente
elegante.

\begin{fichaeixo}{Invariantes de negócio e operação da Plataforma Unooria}
\begin{campos}
  \item[Modelo B2G] \textbf{1 contrato = 1 stack isolada de infraestrutura},
    provisionada por Terraform em projeto/conta própria. Dentro do contrato,
    100+ escolas compartilham a stack com isolamento lógico estrito por
    \cd{escola\_id} (\emph{Row-Level Security} no PostgreSQL).
  \item[Escala por contrato] Cerca de \textbf{\num{10000} estudantes} ativos,
    com picos concentrados de \textbf{\num{2000} a \num{4000} usuários
    simultâneos} no início do turno escolar (8h da manhã) e vales profundos à
    noite e nos fins de semana.
  \item[Equipe de operação] \textbf{Uma pessoa} desenvolve o código, mantém os
    módulos Terraform, responde a incidentes e publica versões. Todo serviço que
    exige manutenção de cluster (como Kubernetes/EKS/AKS/GKE) ou ajustes manuais
    de rede/certificados é descartado em favor de serviços gerenciados
    \emph{serverless}.
  \item[Zero credenciais longas] Nenhuma chave estática de API ou JSON de conta
    de serviço em código, \arq{.env} de produção ou segredos do GitHub. O CI/CD
    autentica por federação OIDC (\emph{Workload Identity}) e os serviços em
    execução usam a identidade nativa do ambiente (ADC / Task Role / Managed
    Identity).
\end{campos}
\end{fichaeixo}

\subsection{Requisitos técnicos determinantes por eixo}

A Tabela~\ref{tab:requisitos-determinantes} sintetiza os requisitos técnicos que
funcionam como \enquote{filtro eliminatório} ao escolher o serviço correspondente
em AWS e Azure.

\begin{table}[H]
\centering
\caption{Requisitos técnicos da Unooria que determinam a escolha dos serviços em nuvem}
\label{tab:requisitos-determinantes}
\begin{tabularx}{\linewidth}{@{}L{28mm}L{44mm}X@{}}
\toprule
\headrow \thd{Eixo} & \thd{Requisito crítico} & \thd{Por que elimina soluções comuns de mercado}\\
\midrule
\textbf{Plataforma} &
  Borda única com roteamento por \emph{hostname}, WAF OWASP e certificados
  validados por DNS sem janela sem TLS &
  O Hub, a API, as Experiências e os Assets respondem sob domínios próprios no
  mesmo balanceador. Adicionar um domínio nunca pode recriar o certificado de
  quem já está no ar (\adrr{ADR-57}, \arq{loadbalancer.tf}).\\
\zebra \textbf{Plataforma} &
  PostgreSQL 16 privado com \cd{FORCE ROW LEVEL SECURITY}, particionamento
  mensal e duas \emph{roles} sem \cd{BYPASSRLS} &
  Descarta bancos NoSQL e exige PostgreSQL gerenciado real, sem IP público, com
  acesso administrativo exclusivamente via túnel autenticado por IAM (IAP/SSM) e
  \emph{connection pooling} para absorver escalonamento de contêineres.\\
\textbf{Plataforma} &
  Fila de telemetria com entrega \emph{push} HTTP autenticada por OIDC e DLQ &
  A API centraliza a ingestão em \rota{POST /telemetry/ingest}. A nuvem deve
  empurrar a mensagem autenticada para a própria API sem exigir um serviço
  \emph{worker} separado rodando 24/7 só para fazer \emph{polling} de fila.\\
\zebra \textbf{Experiências} &
  \emph{Bundles} estáticos imutáveis em \cd{/\{expId\}/\{versao\}/} com origem
  própria no CDN &
  Experiências nunca rodam em contêiner (\emph{regra inviolável}). Precisam de
  domínio dedicado (\cd{exp\_domain}) para que o \cd{targetOrigin} do
  \cd{postMessage} do Hub não compartilhe origem com o restante da nuvem.\\
\textbf{API} &
  Contêiner \emph{serverless} com alta concorrência (80 req/instância), CPU
  sempre alocada e \emph{timeout} de até \num{3600}\,s &
  A geração de planos e pacotes responde \cd{HTTP 202} e segue processando em
  segundo plano no mesmo processo; o Agente do Professor mantém SSE por até
  300\,s; o mascote da Ajuda mantém WebSocket bidirecional por até 60 minutos.\\
\zebra \textbf{Uso de IA} &
  Saída longa (64k \emph{tokens}), controle de orçamento de raciocínio, geração
  de imagem com \emph{aspect ratio} nativo (\cd{4:5}) e voz bidirecional &
  Planos de aula anuais cortam em 32k \emph{tokens}; chamadas rápidas em sala de
  aula (\cd{diagramar}) não podem estourar latência nem orçamento oculto de
  \emph{thinking tokens} (\adrr{ADR-148}).\\
\textbf{Uso de IA} &
  Runtime gerenciado para agente Python (ADK) com sessões persistentes + busca
  estruturada (BNCC) e por \emph{layout} de PDF com \cd{pageSpan} &
  O Agente do Professor precisa listar 100\% das habilidades de um ano por
  filtro exato e citar livro + intervalo exato de páginas (\cd{p. X--Y}) sem
  manter banco vetorial provisionado por hora a custo proibitivo.\\
\bottomrule
\end{tabularx}
\end{table}

\begin{regra}
\textbf{Simplicidade operacional é requisito funcional.} Quando uma nuvem oferece
dois caminhos para o mesmo problema --- por exemplo, rodar a API em Kubernetes
(EKS/AKS) versus rodá-la em um serviço de contêineres \emph{serverless} (ECS
Fargate / Azure Container Apps) ---, este estudo seleciona sempre o caminho
\textbf{gerenciado e sem cluster}, preservando a regra de negócio de que toda a
plataforma é operada por uma única pessoa.
\end{regra}
